在问题二中，我们把上游资源-鱼群的动态轨迹先汇总为综合生态状态指数，用该指数驱动中层猎物承载力，再建立长江江豚、长江鲟及入侵物种的种群响应模型。高阻隔、无人工放流和污染分别作为对照情景。基线情景下，长江江豚与长江鲟种群末值分别为禁渔前的1.133倍和1.140倍；高阻隔情景下分别降至0.901倍和0.856倍；无人工放流时长江鲟降至0.519倍；污染情景下二者进一步降至0.610倍和0.704倍。情景之间的差异说明，江湖通道阻隔同时压低两类珍稀物种的恢复水平，而人工增殖放流对长江鲟的补偿更明显。